\cotefile{58}
\scanpage{2}
\begin{flushright}
\uline{\unsure{Application|Applications|Applicat.} à} $\pi_0$
\end{flushright}
% archive sheet number: 1

\scanpage{3}
\deleted{les composantes de $X$ de dim $\geq 3$}

\uline{Lemme 2} \quad $A$ anneau local \unsure{noeth.}, \inserted{complet}, $f\in\mathfrak rA$ \quad $X=\operatorname{Spec} A$, $X'=X-\{\mathfrak m\}$, $Y'=X'\cap V(f)$.

Supposons \inserted{\unsure{que|pour} \deleted{\unsure{les}} \unsure{ttes|tte}} \deleted{les} composantes \deleted{\unsure{de} $X'_i$ \unsure{de dim} $\geq 2$,} \inserted{de $X'$, \deleted{\unsure{et dim}}} \deleted{les $X_i$ de $X$, les \unsure{composantes} de $\widehat X_i$ sont de dim $\geq 3$} [\unsure{donc si}, $A$ est \unsure{quotient|composé} \deleted{[\unsure{Alors} \unsure{de} $X'_i$, \unsure{les composantes sont de dim} $\geq 1$. $A$.]} d'un régulier, \deleted{il suffit pour ceci que} ceci signifie \unsure{simplem}$^{\text{t}}$ que \deleted{les $X'_i$ sont de dim $\geq 2$}] de dim $\geq 2$. Alors les composantes de $Y'$ sont de dim $\geq 1$, $Y'$ est \uline{connexe}.

En effet, \deleted{l'\unsure{inégalité|intersection} des dim}, \inserted{\unsure{pour} \unsure{tte} partie \inserted{\unsure{fermée, irréduct|fermée-ouverte}} $Z'$ de $Y'$}, les composantes de $Z'\cap Y'$ sont de dim $\geq\dim Z'-1$\inserted{, et $Z'\cap Y'\ne\varnothing$ si $\dim Z'>0$}; \deleted{\unsure{donc|les} si} \inserted{en particulier (si $Z'=X'$)} les composantes de $Z'$ sont de dim $\geq 1$. De plus, pour tte partie $Z'$ de dim $\geq 1$, $Z'\cap Y'$ est de dim $\geq 0$. Donc appliquant le Cor. 1 au Lemme 1, pour prouver le Lemme 2, on est ramené au cas $X'$ \deleted{\unsure{connexe}} irréductible, et si on veut $X$ intègre. de dim $\geq 2$. [D'ailleurs, on \unsure{pourrait|pourra} de prime abord se ramener au cas où $A$ \unsure{était|étant} \uline{complet}, car l'hypothèse porte essentiellement sur le complété (N.B. $\widehat Y$ connexe implique $Y$ connexe, comme image continue de $\widehat Y$).] Alors $\widetilde A$ clôture normale
% archive sheet number: 2

\scanpage{4}
\mspage{- 7 -}
\begin{flushright}
n\textsuperscript{o}\handwritten{354}
\end{flushright}

\uline{est un sous-module gradué} de $M$.

COROLLAIRE 1.- \uline{Si} $N$ \uline{est un sous-module gradué de} $M$ \uline{et si} $(x_{\handwritten{\unsure{\iota|\tau}}})$ \uline{est un système de générateurs de} $N$, \uline{les composantes homogènes des} $x_{\handwritten{\iota}}$ \uline{forment un système de générateurs de} $N$.

COROLLAIRE 2.- \uline{Si} $N$ \uline{est un sous-module gradué de type fini} $M$, $N$ \uline{admet un système fini de générateurs homogènes}.

Il suffit d'appliquer le cor.1 en remarquant qu'un élément de $M$ n'a qu'un nombre fini de composantes homogènes $\ne 0$.

Un sous-module gradué de $A_s$ (resp. $A_d$) est appelé \uline{idéal à gauche} (resp. à \uline{droite}) gradué de l'anneau gradué $A$. Si un sous-anneau $B$ de $A$ est un \uline{sous}-$Z$-module gradué de $A$, on a $(B\cap A_\lambda)(B\cap A_\mu)\subset B\cap A_{\lambda+\mu}$ donc la graduation $(B\cap A_\lambda)$ de $B$ est compatible avec sa structure d'anneau ; on dit alors que $B$ est un \uline{sous-anneau gradué} de $A$.

Il est clair que si $N$ (resp. $B$) est un sous-$A$-module gradué de $M$ (resp. un sous-anneau gradué de $A$), l'injection canonique $N\to M$ (resp. $B\to A$) est un homomorphisme gradué de modules (resp. d'anneaux) de degré 0.

Si $N$ est un sous-module gradué d'un $A$-module gradué $M$, $(M_\lambda)$ la graduation de $M$, les $(M_\lambda+N)/N$, identifiés canoniquement à des sous-modules du module quotient $M/N$, forment une \uline{graduation} compatible avec la structure de ce module pour $\lambda\in\Delta$ ; en effet, ces sous-modules s'identifient aux quotients $M_\lambda/(M_\lambda\cap N)$, et $M/N$ est donc leur somme directe ; en outre, on a $A_\lambda(M_\mu+N)\subset A_\lambda M_\mu+N\subset M_{\lambda+\mu}+N$ donc $A_\lambda((M_\mu+N)/N)\subset(M_{\lambda+\mu}+N)/N$, ce qui établit notre assertion. On dit que la graduation $((M_\lambda+N)/N)$ \uline{est la graduation quotient} de celle de $M$ par $N$, et le module quotient $M/N$, muni de cette graduation, s'appelle le \uline{module quotient gradué} de $M$ par $N$ ; l'homomorphisme canonique $M\to M/N$ est un homomorphisme gradué de degré 0 pour cette graduation.
% archive sheet number: 3

\scanpage{5}
de $A$ est fini sur $A$, et d'ailleurs un anneau local. Appliquant Cor. 2 du lemme, on est ramené au cas $A$ \uline{local} \uline{complet} \inserted{\uline{normal}} de dim $\geq\deleted{\unsure{2}}\unsure{3|2}$, donc $X'$ normal de dim $\geq 2$. Alors
\[
 H^0(\widehat X',\underline{\mathcal O}_{X'})\xleftarrow{\sim}
 H^0(X',\underline{\mathcal O}_{X'})\qquad(\text{\unsure{cf|cif|cf.}})
\]
d'où la conclusion.

\uline{Lemme 3} \quad $A$ anneau local \deleted{\unsure{complet}} \inserted{\unsure{noeth.}} complet. \deleted{\unsure{Supposons} $f$ \unsure{appartenant} \unsure{dans les points} \unsure{de $A$|$\in\mathfrak r A$} \unsure{génériques} \illegible{2}.}

$X$, $X'$, \deleted{[$Y'=V(f)\cap X'$]}. Supposons \inserted{$k\geq 1$ et} que

(i) les composantes de $X'$ de dim $\geq k+1$

(ii) $X'$ est connexe en dim $\geq k$

Alors \inserted{$\exists\ \Phi$ fini de $X'$ tel que $f\in\unsure{\mathfrak r A|\mathfrak r(A)}$, $V(f)\cap\Phi=\varnothing$} \deleted{$\exists f\in\mathfrak rA$ tel que} $\Longrightarrow Y'=X'\cap V(f)$ \deleted{Alors $Y'$ est \unsure{de sorte} \unsure{que} \unsure{de dim|en dim} $\geq k$,} \deleted{\unsure{et|est}} connexe en dim $\geq k-1$.

De façon précise, on construit une partie finie $\Phi$ de $X'$ telle que $V(f)\cap\Phi=\varnothing\Longrightarrow$ conclusion voulue.

Il suffit par le cor 1 : lemme 1 que $V(f)\cap\unsure{X'_i|X'}$ soient connexes en dim $\geq k-1$.
% archive sheet number: 4

\scanpage{6}
\mspage{- 6 -}
\begin{flushright}
n\textsuperscript{o}\handwritten{354}
\end{flushright}

\uline{dués de degré} 0. Mais on notera qu'un homomorphisme gradué bijectif\handwritten{\unsure{'}} de\handwritten{\unsure{''}} degré $\ne 0$ n'est pas un isomorphisme de $A$-modules gradués si $M$ et $N$ ne sont pas réduits à 0.

\uline{Exemples}.- 9) Si $M$ est un $A$-module gradué, $M(\lambda_0)$ un $A$-module gradué obtenu par décalage (Exemple 6), l'application identique de $M(\lambda_0)$ sur $M$ est un homomorphisme gradué bijectif de degré $\lambda_0$.

10) Si $a$ est un élément homogène de degré $\delta$ appartenant au centre de $A$, l'endomorphisme $x\to ax$ d'un $A$-module gradué quelconque $M$ est un homomorphisme gradué de degré $\delta$.

2. \uline{Sous-modules gradués.}

PROPOSITION 2.- \uline{Soit} $\Delta$ \uline{un groupe commutatif}. \uline{Soient} $A$ \uline{un anneau gradué de type} $\Delta$, $M$ \uline{un} $A$-\uline{module gradué de type} $\Delta$, $(M_\lambda)$ \uline{sa graduation}, $N$ \uline{un sous}-$A$-\uline{module de} $M$. \uline{Les propriétés suivantes sont équivalentes} :

a) $N$ \uline{est somme des} $N\cap M_\lambda$ $(\lambda\in\Delta)$.

b) \uline{Les composantes homogènes de tout élément de} $N$ \uline{appartiennent à} $N$.

c) $N$ \uline{est engendré par des éléments homogènes}.

Tout élément de $\unsure{N|M}$ s'écrivant d'une seule manière comme somme d'éléments des $M_\lambda$, il est immédiat que a) et b) sont équivalentes, et que a) entraîne c). Montrons que c) implique b). Soit donc $(x_{\handwritten{\unsure{\iota|\tau}}})_{\handwritten{\iota}\in I}$ une famille de générateurs homogènes $\ne 0$ de $N$ et soit $\delta(\handwritten{\iota})$ le degré de $x_{\handwritten{\iota}}$. Tout élément $x\in N$ s'écrit $\sum\limits_{\handwritten{\iota}}a_{\handwritten{\iota}}x_{\handwritten{\iota}}$ avec $a_{\handwritten{\iota}}\in A$ ; si $a_{\handwritten{\iota},\lambda}$ est la composante de degré $\lambda$ de $a_{\handwritten{\iota}}$ $(\lambda\in\Delta)$, on peut donc écrire $x=\sum\limits_{\lambda\in\Delta}\bigl(\sum\limits_{\handwritten{\iota}\in I}a_{\handwritten{\iota},\lambda-\delta(\handwritten{\iota})}x_{\handwritten{\iota}}\bigr)$ ; comme $\sum\limits_{\handwritten{\iota}\in I}a_{\handwritten{\iota},\lambda-\delta(\handwritten{\iota})}x_{\handwritten{\iota}}$ est homogène de degré $\lambda$ et appartient à $N$, on voit que les composantes homogènes de $x$ appartiennent à $N$.

Lorsqu'un sous-module $N$ de $M$ vérifie les propriétés équivalentes énoncées dans la prop.2, il est clair que les $N\cap M_\lambda$ forment une \uline{graduation} compatible avec la structure de module de $N$ ; on dit que $N$, muni de cette graduation,
% archive sheet number: 5

\scanpage{7}
Cela nous ramène au cas de $X$ irréductible. Par le corollaire 2 au lemme 1, on est ramené au cas où $X$ est de plus \uline{normal}, et à fortiori $S_2$. Or \uline{on \unsure{prouve} dans ce cas qu'il existe une partie} \uline{finie} $\Phi$ de $X'$ telle que $V(f)\cap\Phi=\varnothing$ implique que $V(f)\cap X'=Y'$ est $(S_2)$, [en interprétant $(S_2)$ en \unsure{terme|termes} de codim des $Z'_i$ dans $X'$ ; on \unsure{prend} pour $\Phi$ \unsure{l'ens} des \unsure{pts} minimaux des \deleted{composantes} $Z'_i$].

Mais alors (Hartshorne) $Y'$ est \uline{connexe en codim 1}, donc comme ses composantes \unsure{sont} de dim $\geq k$, $Y'$ est connexe en dim $\geq k-1$.
% archive sheet number: 6

\scanpage{8}
\mspage{- 5 -}
\begin{flushright}
n\textsuperscript{o}\handwritten{354}
\end{flushright}

2) Les $\rlap{\ensuremath{\handwritten{\underline{\underline{\phantom{Z}}}}}}Z$-modules gradués de type $\Delta$, lorsque $\rlap{\ensuremath{\handwritten{\underline{\underline{\phantom{Z}}}}}}Z$ est muni de la graduation triviale, ne sont autres que les groupes commutatifs gradués de la déf. 1.

DEFINITION 4.- \uline{Soient} $A$, $A'$ \uline{deux anneaux gradués de même type} $\Delta$, $(A_\lambda)$, $(A'_\lambda)$ \uline{leurs graduations respectives}. \uline{On dit qu'un homomorphisme d'anneaux} $h:A\to A'$ \uline{est gradué si l'on a} $h(A_\lambda)\subset A'_\lambda$ \uline{pour tout} $\lambda\in\Delta$. \uline{Soient} $M$, $M'$ \uline{deux modules gradués de type} $\Delta$ \uline{sur un anneau gradué} $A$ \uline{de type} $\Delta$, $(M_\lambda)$, $(M'_\lambda)$ \uline{les graduations respectives de} $M$ \uline{et} $M'$. \uline{On dit qu'un} $A$-\uline{homomorphisme} $u:M\to M'$ \uline{est gradué de degré} $\delta\in\Delta$ \uline{si l'on a} $u(M_\lambda)\subset M'_{\lambda+\delta}$ \uline{pour tout} $\lambda\in\Delta$.

Si $u\ne 0$ et si tout élément de $\Delta$ est \uline{régulier}, le degré $\delta$ de $u$ est déterminé de façon \uline{unique}. Si $h:A\to A'$, $h':A'\to A''$ sont deux homomorphismes gradués d'anneaux gradués de type $\Delta$, il en est de même de $h'\circ h:A\to A''$ ; pour qu'une application $h:A\to A'$ soit un isomorphisme d'anneaux gradués, il faut et il suffit que $h$ soit bijective et que $h$ et l'application récipro- $h'$ soient des homomorphismes gradués ; il suffit d'ailleurs pour cela que $h$ soit un homomorphisme gradué bijectif. On voit donc qu'on peut prendre pour \uline{morphismes} de l'espèce de structure d'anneau gradué de type $\Delta$ les homomorphismes gradués (\uline{Ens}., Chap.IV, \S{} 2,n\textsuperscript{o}1).

De même, si $u:M\to M'$ et $v:M'\to M''$ sont deux homomorphismes gradués $\ne 0$ de $A$-modules gradués de type $\Delta$, de degrés respectifs $\delta$ et $\delta'$, $u'\handwritten{\unsure{\circ|\mathrm{o}}}u:M\to M''$ est un homomorphisme gradué de degré $\delta+\delta'$. Si $\delta$ est inversible et si $u:M\to M'$ est un homomorphisme gradué bijectif de degré $\delta$, l'application réciproque $u':M'\to M$ est un homomorphisme gradué bijectif de degré $-\delta$. On en conclut comme ci-dessus qu'on peut prendre pour \uline{morphismes} de l'espèce de structure de $A$-module gradué de type $\Delta$, les \uline{homomorphismes gra}-
% archive sheet number: 7

\scanpage{9}
\fbox{\uline{Théorème}} \quad $A$ anneau local noethérien, \inserted{$X$, $X'$, $k\geq 1$}

\tikz[remember picture,overlay]{\coordinate (Ktheoremstart) at (-.6em,1.55ex);}\deleted{tel que} \inserted{Supposons $a_k$)} les composantes irréductibles de $\widehat X=\operatorname{Spec}(\widehat A)$ soient de dim $\geq k+2$

% slanted
\margin{\uline{\unsure{Canulé|annulé|Conditions}}, faire \uline{ttes} les hypothèses sur $\widehat A$, pour en déduire des conclusions sur $A$ [la réduction au cas complet étant alors comme ici]}
[il suffit pour ceci que celles de $\unsure{X|N}=\operatorname{Spec} A$ soient de dim $\geq k+2$ i.e. celles de $X'$ de dim $\geq k+1$, et cela est suffisant si p. ex. $A$ est quotient d'un régulier,]. \deleted{\unsure{Soient} $f_1$} $b_k$) $X'$ est connexe en dim $\geq k$.

Considérons \deleted{des \unsure{éléments}} $0\leq\ell\leq k$, et $f_1\mathinner{\text{---}}f_\ell\in\mathfrak rA$, \deleted{$\unsure{H_i|Y_i}=V(f_i)$}, $Y'=X'\cap\bigcap\limits_i V(f_i)$. Alors \deleted{les composantes de} $Y'$ satisfait les conditions $a_{k-\ell}$, $b_{k-\ell}$, notamment $Y'$ est connexe en dim $\geq k-\ell$..\tikz[remember picture,overlay]{\coordinate (Ktheoremend) at (0,-.4ex); \draw[thin] (Ktheoremstart) -- (Ktheoremstart |- Ktheoremend);}

Pour $a_{k-\ell}$ c'est immédiat. Reste $b_{k-\ell}$. Si on pouvait \unsure{desconnecter|désconnecter|déconnecter} $Y'$ par $Z'$ de dim $<k-\ell$, il \unsure{existerait|existera|existence} \inserted{$k-\ell$ \unsure{élts|élémts} de $\mathfrak m A$,} $f_{\ell+1}\unsure{-|\ldots}f_k$ tels que $Z'\cap\deleted{\unsure{X'}}V(f_{\ell+1})\cap-\cap V(f_k)=\varnothing$, d'où \unsure{résulterait}
% archive sheet number: 8

\scanpage{10}
\mspage{- 4 -}
\begin{flushright}
n\textsuperscript{o}\handwritten{354}
\end{flushright}

5) Soient $A$ un anneau gradué de type $\Delta$, $M$ un $A$-module gradué de type $\Delta$, $\handwritten{\rho}$ un homomorphisme de $\Delta$ dans un monoïde commutatif $\Delta'$ ayant un élément neutre, tel que $\handwritten{\rho}(0)=0$. Alors $A$ est un anneau gradué de type $\Delta'$ et $M$ un $A$-module gradué de type $\Delta'$ pour les graduations de type $\Delta'$ déduites des graduations de type $\Delta$ de $A$ et $M$ par le procédé de l'exemple 2 ; cela résulte aussitôt de la relation $\handwritten{\rho}(\lambda+\lambda')=\handwritten{\rho}(\lambda)+\handwritten{\rho}(\lambda')$.

En particulier, si $\Delta=\Delta_0\times\Delta_0$, on dira que $A$ (resp. $M$) est un anneau (resp. un $A$-module) \uline{bigradué} de type $\Delta_0$ ; les graduations partielles et la graduation totale correspondantes, de type $\Delta_0$, sont alors compatibles avec la structure d'anneau de $A$ (resp. la structure de module de $M$).

6) Soient $A$ un anneau gradué de type $\Delta$, $M$ un $A$-module gradué de type $\Delta$, $\lambda_0$ un élément de $\Delta$ ; pour tout $\lambda\in\Delta$, posons $M'_\lambda=M_{\lambda+\lambda_0}$ ; comme $A_\lambda M'_\mu\subset M_{\lambda+\mu+\lambda_0}=M'_{\lambda+\mu}$, les $M'_\lambda$ forment sur $M$ une nouvelle graduation de type $\Delta$ ; le $A$-module gradué de type $\Delta$ ainsi obtenu est noté $M(\lambda_0)$ et on dit qu'il s'obtient par \uline{décalage} de $\lambda_0$ de la graduation de $M$.

*7) Soit $B$ un anneau commutatif. L'anneau de polynômes $B[X]$ à une indéterminée est gradué de type $\rlap{\ensuremath{\handwritten{\underline{\underline{\phantom{N}}}}}}N$ par les sous-groupes $BX^n$ $(n\geq 0)$ (cf. Chap. III et IV).

8) Soient $B$ un anneau commutatif, $E$ un $B$-module, $Q$ une forme quadratique sur $E$, $C(Q)$ l'algèbre de Clifford de $Q$ (cf. Chap.IX,\S{}9) ; les sous-$B$-modules $C^+(Q)$ et $C^-(Q)$ forment sur $C(Q)$ une graduation de type $\rlap{\ensuremath{\handwritten{\underline{\underline{\phantom{Z}}}}}}Z/2\rlap{\ensuremath{\handwritten{\underline{\underline{\phantom{Z}}}}}}Z$ compatible avec la structure d'anneau de $C(Q)$.*

\uline{Remarques.}- 1) Les graduations les plus souvent utilisées sont de type $\rlap{\ensuremath{\handwritten{\underline{\underline{\phantom{Z}}}}}}Z$ ou de type $\rlap{\ensuremath{\handwritten{\underline{\underline{\phantom{Z}}}}}}Z^n$ ; lorsqu'on parle d'anneaux ou de modules \uline{gradués} (resp. \uline{bigradués}, \uline{trigradués}, etc.) sans préciser de quel type, il est sous-entendu qu'il s'agit de graduations de type $\rlap{\ensuremath{\handwritten{\underline{\underline{\phantom{Z}}}}}}Z$ (resp. $\rlap{\ensuremath{\handwritten{\underline{\underline{\phantom{Z}}}}}}Z^2$, $\rlap{\ensuremath{\handwritten{\underline{\underline{\phantom{Z}}}}}}Z^3$, etc.) ; on dit aussi qu'un anneau (resp. module) gradué de type $\rlap{\ensuremath{\handwritten{\underline{\underline{\phantom{N}}}}}}N$ est un anneau (resp. module) \uline{gradué à degrés positifs}.
% archive sheet number: 9

